(3 points + 3 points)
\begin{teilaufgaben}
\item
The function $u(x)$ is defined on the interval $\Omega = [-1, 1].$

The function satisfies in $\Omega$ 
\[
u''(x) = x^2
\]
and $u(-1) = u(1) = 1$ on the boundary of $\Omega$.

Determine an approximation function $\tilde u(x)$ for $u(x)$ by
the method of Galerkin. 

Use the ansatz $\tilde u(x) = 1 + a_1 \cdot (1-x^2) + a_2 \cdot (1-x^4).$

\item
The function $u(x)$ is defined on the interval $\Omega = [0, 1].$

The function satisfies in $\Omega$ 
\[
u''(x) = - x
\]
and $u(0) = 0$ and $u(1) = 0$ on the boundary of $\Omega$.  

Determine an approximate value for $u(1/2)$ by the method of finite elements.

Use quadratic finite elements with $h = 1$.
\end{teilaufgaben}

\begin{loesung}
\begin{teilaufgaben}
\item
Plugging the ansatz into the differential equation
\[
u''(x) - x^2 = 0
\]
we obtain
\begin{equation}
a_1 \cdot (-2) + a_2 \cdot (-12 \cdot x^2) - x^2 \approx 0.
\tag{{\bf 1P}}
\end{equation}
Since
\[
\int_{-1}^1 1 - x^2 \ dx = 4/3,
\qquad
\int_{-1}^1 (1-x^2) \cdot x^2 dx = 4/15,
\]
as well as
\[
\int_{-1}^1 1 - x^4 \ dx = 8/5,
\qquad
\int_{-1}^1 (1 - x^4) \cdot x^2 \ dx = 8/21,
\]
we obtain by Galerkin the linear system of equations
\begin{equation}
\left[\begin{array}{cc}
  -8/3 & -16/5 \\
 -16/5 & -32/7
\end{array}\right]
\cdot
\left(\begin{array}{r} a_1 \\ a_2 \end{array}\right)
-
\left(\begin{array}{c} 4/15 \\ 8/21 \end{array}\right)
=
\left(\begin{array}{r} 0 \\ 0 \end{array}\right).
\tag{{\bf 1P}}
\end{equation}
Therefore $(a_1, a_2) = (0, -1/12)$ and
\begin{equation}
\tilde u(x) =  1/12 \cdot (x^4 + 11).
\tag{{\bf 1P}}
\end{equation}
\item
The corresponding Ritz matrix is
\begin{equation}
R = 1/3 \cdot
\left[\begin{array}{rrr}
 7 & -8 &  1 \\
-8 & 16 & -8 \\
 1 & -8 &  7
\end{array}\right],
\tag{{\bf 1P}}
\end{equation}
whereas the Ritz vector is
\[
\underline{r} = (*,r,*)
\]
with
\begin{equation}
r = \int_0^1 4 \cdot (x - x^2) \cdot x \ dx = 1/3.
\tag{{\bf 1P}}
\end{equation}
The vector of the knot variable is
\[
\underline{a} = (0, a, 0).
\]
Reduction of the Ritz system leads to  
\[
16/3 \cdot a = 1/3.
\]
Therefore we have
\begin{equation}
\tilde u(1/2) = 1/16
\tag{{\bf 1P}}
\end{equation}
\end{teilaufgaben}
\end{loesung}
